\documentclass[lettersize,journal]{IEEEtran}
\usepackage[T1]{fontenc}
\usepackage{amsmath}
\usepackage{url}
\usepackage{cite}

\begin{document}

\title{CS 550 Final Project Proposal:\\Replication of the Hawkeye Cache
Replacement Policy}

\author{Vito Wang, Endong Dai, and Xiaotian Tong}

\maketitle

\begin{abstract}
This project replicates the Hawkeye cache replacement policy~\cite{hawkeye}
in gem5 and evaluates it on SPEC CPU2017. Hawkeye learns from Belady's optimal
replacement algorithm. It reconstructs the decisions that the optimal policy
would have made on past cache accesses and uses them to train a predictor that
decides which cache lines to keep. We will validate each component with unit
tests and logs, check that programs produce the same results as under LRU, and
compare Hawkeye's miss rate and performance with the replacement policies that
gem5 already provides, such as LRU, DRRIP, and SHiP.
\end{abstract}

%% ------------------------------------------------------------------
\section{Introduction \& Problem Statement}

The last-level cache (LLC) is the last on-chip structure that can hide the long
latency of DRAM, so every LLC miss stalls the processor for hundreds of cycles.
When the LLC is full, the replacement policy chooses which line to evict, and
this choice largely decides the miss rate. The common least-recently-used (LRU)
policy works well when recently used data is soon reused. However, it fails when
a program cycles through a working set larger than the cache, because every
line is then evicted just before it is reused, and when streaming data that is
never reused pushes useful lines out. Improved policies handle some of these
patterns, for example by giving new lines a low priority so that they leave
the cache quickly unless they are reused (RRIP~\cite{rrip}), or by learning
which load instructions bring in lines that are reused (SHiP~\cite{ship}).
However, each is built around a particular class of access patterns.

Belady's algorithm~\cite{belady}, also called OPT, is the optimal policy. It
evicts the line whose next use lies furthest in the future. OPT cannot be built
in hardware because it needs future knowledge. Hawkeye~\cite{hawkeye} avoids
this problem by looking at the past instead. It computes what OPT would have
done on past accesses, learns which load instructions tend to bring in lines
that OPT would keep, and uses this to guide future evictions. The original paper
evaluated Hawkeye in the Cache Replacement Championship simulator on SPEC CPU2006.
gem5 does not include Hawkeye, so it has not been evaluated with gem5's
out-of-order CPU model or with SPEC CPU2017.

This project aims to implement Hawkeye~\cite{hawkeye} in
gem5~\cite{gem5,gem5v20}. The goal is to verify its functional correctness and
to compare its miss rate and performance with the baseline policies in gem5.

%% ------------------------------------------------------------------
\section{Related Work}

Classic replacement policies rank lines by recency or frequency, e.g., LRU and
LFU~\cite{hp6}. RRIP~\cite{rrip} instead predicts how soon each line will be
re-referenced, using a 2- or 3-bit counter per line, and inserts new lines with
a long predicted interval so that data that is never reused leaves the cache
quickly. DRRIP chooses between two RRIP variants at run time through set
dueling~\cite{dip}, which compares the policies on a few dedicated sets. A
previous CS 550 project evaluated DRRIP, SRRIP, and BRRIP against LRU in gem5.
We use these policies as baselines.

Other policies predict reuse for each load instruction. SHiP~\cite{ship} learns,
for each program counter (PC), whether the lines it inserts are reused, and
SDBP~\cite{sdbp} predicts dead blocks from a small sampler. Both learn from the
behavior of a heuristic policy, so they inherit its limits.

Hawkeye~\cite{hawkeye} instead learns from OPT itself. Its authors won the
2nd Cache Replacement Championship (CRC2)~\cite{crc2}, and we will use the CRC2
code as a reference. Its follow-up, Mockingjay~\cite{mockingjay}, predicts reuse
distances instead of a binary decision. We leave Mockingjay out of scope.

%% ------------------------------------------------------------------
\section{Research Questions}

\begin{itemize}
\item \textbf{What is the Hawkeye policy, and can its reconstruction of OPT be
implemented correctly in gem5?}
\item \textbf{How do Hawkeye's LLC miss rate and IPC compare with LRU, DRRIP,
and SHiP?}
\item \textbf{Why can Hawkeye reduce misses in access patterns where LRU and
RRIP fail?}
\end{itemize}

%% ------------------------------------------------------------------
\section{Proposed Solution}

Jain and Lin proposed Hawkeye~\cite{hawkeye}, which has three parts.

\textbf{OPTgen} analyzes past accesses and reconstructs the decisions that OPT
would have made. When a line is accessed again, OPTgen checks whether OPT could
have kept the line in the cache since its previous access without exceeding the
cache capacity. If so, OPT would have hit. Otherwise, OPT would have missed. To
save storage, OPTgen observes only a small sample of the cache sets.

\textbf{The Hawkeye Predictor} learns from OPTgen which load instructions bring
in lines that are worth keeping. It keeps a small counter for each program
counter (PC). When OPTgen decides that OPT would have hit, the counter of the
PC that last accessed the line goes up. Otherwise, it goes down. When a line is
accessed, the counter of its PC predicts whether the line is
\emph{cache-friendly} or \emph{cache-averse}.

\textbf{Replacement} decides which line to evict based on these predictions.
Cache-averse lines are evicted first. If a set has none, Hawkeye evicts the
oldest cache-friendly line and lowers the counter of the PC that brought it
in.

%% ------------------------------------------------------------------
\section{Evaluation}

\subsection{Methodology}

We will use gem5 v23.0.0.1 on the course AWS instance, in syscall-emulation (SE)
mode with the x86 out-of-order CPU model (O3CPU). Following the original
paper, the cache hierarchy has 32\,KB 4-way L1 instruction and data caches, a
256\,KB 8-way L2, and a 2\,MB 16-way LLC. The course configuration has only two
cache levels, so we will add the LLC. We will add the following components to
gem5:

\begin{itemize}
\item \textbf{Hawkeye replacement policy}: a new class for gem5's
replacement-policy interface, which passes the requesting packet, and thus the
PC, on every insertion and hit, as the existing SHiP policy uses.
\item \textbf{Per-line state}: the age, the PC, and the prediction, stored
with each line because the eviction call does not see the packet.
\item \textbf{OPTgen}: the access history of the sampled sets and the
reconstructed OPT decisions.
\item \textbf{Hawkeye Predictor}: the table of per-PC counters.
\item \textbf{Requests without a PC}: writebacks and other requests from the
L2, which must not train the predictor.
\item \textbf{Logs and statistics}: a debug flag that prints each decision.
\end{itemize}

\subsection{Functional Validation}

The functional correctness of the Hawkeye policy will be tested in two ways.

\emph{a) Unit testing with logs:} We will add logs to each component. For each
access, the log will show OPTgen's hit or miss decision, the predictor counter,
and the age of each line. We will compare them with values calculated by hand for
short access sequences, such as the example in the original paper, and with
values computed by a small script for longer ones.

\emph{b) Execution correctness:} Hawkeye will be integrated into the LLC and
will run a series of test programs. The replacement policy must not change
program results, so the outputs are expected to match those under LRU. Some
test programs will have known access patterns, such as an array that fits in
the LLC and one slightly larger than the LLC, so that we can also check whether
the miss rates match our expectations.

\subsection{Performance Impact}

Hawkeye will be compared with baseline policies that gem5 already provides,
such as LRU, DRRIP, and SHiP. We will focus on memory-intensive SPEC
CPU2017~\cite{spec17} benchmarks, where the LLC matters most, such as mcf,
omnetpp, xalancbmk, and gcc. The comparison will focus on LLC misses per
kilo-instruction (MPKI) and IPC, and we will repeat it for several LLC sizes. We
will also estimate the hardware cost from the storage of the added
structures.

%% ------------------------------------------------------------------
\section{Milestones}

The project is broken down into five milestones that span 5 weeks of work.

\begin{itemize}
\item \textbf{Milestone 1: Setup and Initial Research (Week 1):} We will add
the LLC and the O3CPU to the gem5 configuration, run the baseline policies, study
the Hawkeye paper and the CRC2 code, and confirm that the PC reaches the LLC.
\item \textbf{Milestone 2: Hawkeye Implementation (Week 2):} We will implement
OPTgen, the predictor, and the replacement policy.
\item \textbf{Milestone 3: Functional Validation (Week 3):} We will run the unit
tests and test programs and fix any differences from the expected behavior.
\item \textbf{Milestone 4: Performance Evaluation (Week 4):} We will run SPEC
CPU2017 with every policy and repeat the runs for several LLC sizes.
\item \textbf{Milestone 5: Final Evaluation and Reporting (Week 5):} We will
analyze the results and write the final report and presentation.
\end{itemize}

%% ------------------------------------------------------------------
\section{Division of Labor}

Vito will lead the design and implementation of Hawkeye in gem5, including
OPTgen, the predictor, and the replacement policy. Endong
will design the validation plan, write the test programs and the script that
computes the expected OPTgen decisions, and run the unit tests. Xiaotian will
run the baseline and SPEC CPU2017 simulations, repeat them for several LLC
sizes, and collect and plot the statistics. All three will analyze the results and
write the final report.

%% ------------------------------------------------------------------
\begin{thebibliography}{99}

\bibitem{hawkeye}
A. Jain and C. Lin, ``Back to the future: Leveraging Belady's algorithm for
improved cache replacement,'' in \emph{Proc. ISCA}, 2016.

\bibitem{rrip}
A. Jaleel, K. B. Theobald, S. C. Steely Jr., and J. Emer, ``High performance
cache replacement using re-reference interval prediction (RRIP),'' in
\emph{Proc. ISCA}, 2010.

\bibitem{ship}
C.-J. Wu, A. Jaleel, W. Hasenplaugh, M. Martonosi, S. C. Steely Jr., and
J. Emer, ``SHiP: Signature-based hit predictor for high performance caching,''
in \emph{Proc. MICRO}, 2011.

\bibitem{belady}
L. A. Belady, ``A study of replacement algorithms for a virtual-storage
computer,'' \emph{IBM Systems Journal}, vol. 5, no. 2, pp. 78--101, 1966.

\bibitem{gem5}
N. Binkert \emph{et al.}, ``The gem5 simulator,'' \emph{ACM SIGARCH Comput.
Archit. News}, vol. 39, no. 2, pp. 1--7, 2011.

\bibitem{gem5v20}
J. Lowe-Power \emph{et al.}, ``The gem5 simulator: Version 20.0+,''
arXiv:2007.03152, 2020.

\bibitem{hp6}
J. L. Hennessy and D. A. Patterson, \emph{Computer Architecture: A Quantitative
Approach}, 6th ed. Morgan Kaufmann, 2019.

\bibitem{dip}
M. K. Qureshi, A. Jaleel, Y. N. Patt, S. C. Steely Jr., and J. Emer, ``Adaptive
insertion policies for high performance caching,'' in \emph{Proc. ISCA}, 2007.

\bibitem{sdbp}
S. M. Khan, Y. Tian, and D. A. Jim\'enez, ``Sampling dead block prediction for
last-level caches,'' in \emph{Proc. MICRO}, 2010.

\bibitem{crc2}
The 2nd Cache Replacement Championship (CRC2), held with ISCA 2017.

\bibitem{mockingjay}
I. Shah, A. Jain, and C. Lin, ``Effective mimicry of Belady's MIN policy,'' in
\emph{Proc. HPCA}, 2022.

\bibitem{spec17}
Standard Performance Evaluation Corporation, ``SPEC CPU2017,''
\url{https://www.spec.org/cpu2017/}.

\end{thebibliography}

\end{document}
